\documentclass[10pt]{article}
\usepackage[margin=1in, paperwidth=8.5in, paperheight=11in]{geometry}
\usepackage{ifpdf,amsmath, amssymb, comment, color, graphicx, stmaryrd,setspace,enumitem,tikz, fancyhdr, wrapfig, textcomp, units, mathptmx, pdfpages, cancel, multicol}
\usepackage{hyperref}
\hypersetup{
    colorlinks=true,
    urlcolor=blue,
}

\setlength{\headheight}{14.5pt}
\newcommand{\Q}{\mathbb{Q}}
\newcommand{\R}{\mathbb{R}}
\newcommand{\Z}{\mathbb{Z}}
\newcommand{\vu}{\mathbf{u}}
\newcommand{\vv}{\mathbf{v}}
\newcommand{\vw}{\mathbf{w}}
\newcommand{\vi}{\mathbf{i}}
\newcommand{\vj}{\mathbf{j}}
\newcommand{\vk}{\mathbf{k}}
\newcommand{\vn}{\mathbf{n}}
\newcommand{\vr}{\mathbf{r}}
\newcommand{\va}{\mathbf{a}}
\newcommand{\vF}{\mathbf{F}}
\newcommand{\vL}{\mathbf{L}}
\newcommand{\vT}{\mathbf{T}}
\newcommand{\vN}{\mathbf{N}}
\newcommand{\proj}{\operatorname{proj}}
\newcommand{\orth}{\operatorname{orth}}
\newcommand\del\nabla
\newcommand\dotp[1][.5]{\,\mathbin{\vcenter{\hbox{\scalebox{#1}{$\bullet$}}}}\,}

% Solution text is in red. If you want the solutions to show, remove the \iffalse from the definition of the \red command.
%\newcommand{\red}[1]{ %\iffalse
%	\textcolor{red}{#1} }%\fi}

\newenvironment{red}{\color{red}}{\ignorespacesafterend}
\newenvironment{blue}{\color{blue}}{\ignorespacesafterend}
\newcommand{\green}[1]{\textcolor{green}{#1}}
\renewcommand{\section}[1]{\begin{center} \textbf{#1} \\\end{center}}
%
\hyphenpenalty=5000
\setlength{\parindent}{0in}
%\oddsidemargin=-.25in
\allowdisplaybreaks
\pagestyle{fancy}
\renewcommand{\headrulewidth}{0pt}
\lhead{MATH 202}
\rhead{Spring 2023}
%\lfoot{\copyright\ CLEAR Calculus 2010}
\cfoot{}

\begin{document}
%


%\onehalfspacing
\allowdisplaybreaks
%##################################################################
\section{PS\#4: Integration by parts - \red{Answer key} }

\everymath{\displaystyle}
\begin{enumerate}[leftmargin=0pt]

\item (\href{https://activecalculus.org/single/sec-5-4-parts.html#CBY}{Section 5.4 Exercise 7})
For each of the following indefinite integrals, determine whether you would use 
$u$-substitution, integration by parts, neither, or both to evaluate the integral. In each case, write one sentence to explain your reasoning, and include a statement of any substitutions used. Finally, use your chosen approach to evaluate each integral.
\begin{enumerate}
    \item $\int x^2 \cos(x^3)\ dx$ \begin{red}
        $= \frac13 \sin(x^3) +C$
    \end{red}
    
    \item $\int x^5 \cos(x^3)\ dx = \int x^3\cdot (x^2\cos(x^3))\ dx$
    \begin{red}
        $= \frac13 x^3 \sin(x^3) + \frac13 \cos(x^3) + C$
    \end{red}
    
    \item $\int x\ln(x^2)\ dx$
    \begin{red}
        $= \frac12\left(x^2 \ln(x^2) - x^2\right) + C$
    \end{red}
    
    \item $\int \sin(x^4)\ dx$
    \begin{red}
        : No elementary antiderivative!
    \end{red}
    
    
    \item $\int x^3\sin(x^4)\ dx$
    \begin{red}
        $=-\frac14\cos(x^4)+ C$
    \end{red}
    
    \item $\int x^7 \sin(x^4)\ dx = \int x^4\cdot (x^3\sin(x^4))\ dx$
    \begin{red}
        $=-\frac14 x^4\cos(x^4) + \frac14 \sin(x^4) + C$
    \end{red}
    
\end{enumerate}
\end{enumerate}
\end{document}